% Generated draft: requires mathematical and source review.
\documentclass[11pt]{article}
\usepackage[margin=25mm]{geometry}
\usepackage{fontspec}
\setmainfont{Latin Modern Roman}
\usepackage{amsmath,amssymb,mathtools,unicode-math}
\setmathfont{Latin Modern Math}
\usepackage{enumitem,needspace,xurl,hyperref}
\setlength{\parindent}{0pt}
\setlength{\parskip}{5pt}
\setlength{\emergencystretch}{3em}
\newcommand{\N}{\mathbb{N}}
\newcommand{\Z}{\mathbb{Z}}
\newcommand{\Q}{\mathbb{Q}}
\newcommand{\R}{\mathbb{R}}
\newcommand{\C}{\mathbb{C}}
\newcommand{\F}{\mathbb{F}}
\newcommand{\A}{\mathbb{A}}
\newcommand{\PP}{\mathbb{P}}
\newcommand{\E}{\mathbb{E}}
\newcommand{\eps}{\varepsilon}
\newcommand{\dd}{\,\mathrm{d}}
\begin{document}
\section*{UnsolvedMath ID: 30006416}
\textbf{Source problem code:} catalogue entry
\subsection{Definitions and mathematical statement}
A \(G_2\)-structure on a smooth seven-dimensional manifold \(M\) is a positive differential \(3\)-form \(\varphi\), whose pointwise stabilizer is the compact exceptional group \(G_2\). It determines a Riemannian metric \(g_\varphi\). This terminology alone does not impose a torsion-free condition.

On a Riemannian manifold, a calibration of degree \(k\) is a closed differential \(k\)-form \(\alpha\) satisfying \(\alpha(\xi)\leq 1\) for every oriented unit simple \(k\)-vector \(\xi\). An oriented \(k\)-dimensional submanifold \(Y\) is calibrated by \(\alpha\) when \(\alpha|_Y=\operatorname{vol}_Y\). The supplied record does not specify which calibration or which calibrated submanifolds are intended.

\textbf{Source problem, with unresolved scope.} Define well-defined counting invariants for the intended gauge-theoretic objects and the intended calibrated submanifolds in \(G_2\) geometry, and prove an equality between those invariants.

Schematically, if \(\mathcal D\) denotes the geometric and topological input, including any necessary auxiliary choices, the requested construction would supply invariants
\[
I_{\mathrm{gauge}}(\mathcal D),\qquad I_{\mathrm{cal}}(\mathcal D)
\]
in a common target and establish
\[
I_{\mathrm{gauge}}(\mathcal D)=I_{\mathrm{cal}}(\mathcal D)
\]
for the intended admissible inputs. This display records the requested objective only: the source does not define \(\mathcal D\), either invariant, their common target, or the range of inputs over which equality is sought.

The record does not specify the ambient manifolds, compactness or boundary conditions, torsion assumptions, gauge group, bundles, gauge equations, equivalence relation, calibration, dimensions or topological classes of submanifolds, or counting conventions. It also does not specify the choices under which the counts should be invariant, or how the input data on the two sides should correspond. Consequently, the supplied passage expresses a research programme rather than a fully quantified conjecture; these missing specifications must be resolved before a precise mathematical assertion can be stated.
\subsection{Short English statement}
The problem asks for counting invariants arising from gauge theory and calibrated submanifolds in \(G_2\) geometry, together with an equality relating them. The supplied passage leaves the objects, hypotheses, counting rules, and correspondence unspecified, so the displayed equality is schematic. The supplied open-status assessment is dated curation and requires independent verification.
\subsection{Sources}
\begin{itemize}
\item Local UnsolvedMath registry, permanent ID \(30006416\), title: \emph{Gauge-Theoretic and Calibrated-Submanifold Counting Invariants}; source problem code \texttt{catalogue entry}. The record attributes the problem to an Oberwolfach workshop discussion and names no individual proposer.
\item Supplied report citation: \emph{Differentialgeometrie im Großen (2026)}, Oberwolfach Reports, DOI \texttt{10.4171/owr/2025/30}, pp.\ 1591--1646. \url{https://doi.org/10.4171/owr/2025/30}. The literature-source list labels this reference \emph{Oberwolfach Report 30/2025}. The citation's parenthetical year 2026 differs from the report year 2025 indicated by that label, the DOI, and the supplied proposed year; this discrepancy is unresolved. No precise page or speaker attribution for the passage is supplied.
\item Supplied literature source: \emph{DFG Gauge-theoretic methods in G2 geometry}. \url{https://www.spp2026.de/projects/2nd-period/gauge-theoretic-methods-in-the-geometry-of-g2-manifolds.html}.
\item Literature assessment dated exactly \texttt{2026-08-22}: the supplied curation labels the problem open and states that constructing the two counting invariants and an equality between them remains an explicitly open problem in \(G_2\) geometry. This is supplied dated curation derived from the OWR record, not an independently verified literature assessment; time-sensitive status and cited claims require independent verification.
\item The registry labels the statement \texttt{corrected\_verified}, with the qualification that only formula delimiters, notation, grammar, and typography were repaired while preserving mathematical content and scope. This editorial label does not establish mathematical correctness or resolve the missing hypotheses.
\end{itemize}
\end{document}
